\subsection{FUNDAMENTAL GROUP}

In the following proposition, \(f(\mathrm{G},\mathrm{T})\) denotes the homomorphism from \(\varGamma(\mathrm{T})\) to \(\pi_{1}(\mathrm{G})\) that is the composite of the canonical isomorphism from \(\varGamma(\mathrm{T})\) to \(\pi_{1}(\mathrm{T})\) (no. 2, Remark 3) and the homomorphism \(\pi_{1}(\iota)\) , where \(\iota\) is the canonical injection \(T\to G\) .

PROPOSITION 11. The homomorphism \(f(\mathrm{G},\mathrm{T}) : \Gamma(\mathrm{T}) \to \pi_1(\mathrm{G})\) is surjective. Its kernel is the subgroup \(\mathrm{N}(\mathrm{G},\mathrm{T})\) of \(\Gamma(\mathrm{T})\) generated by the family of nodal vectors \((K_{\alpha})_{\alpha \in \mathrm{R}(\mathrm{G},\mathrm{T})}\) .

The homomorphism \(f(\mathrm{G},\mathrm{T})\) is surjective by Prop. 3 ( \(\S2\) , no. 4). We denote by \(A(\mathrm{G},\mathrm{T})\) the assertion: ``the kernel of \(f(\mathrm{G},\mathrm{T})\) is generated by the \(K_{\alpha}\) '' which it remains to prove, and distinguish several cases:

\begin{enumerate}
\def\labelenumi{\alph{enumi})}
\item
  G is simply-connected. Let \(\rho : \mathfrak{g}_{\mathbf{C}} \to \mathfrak{gl}(\mathrm{V})\) be a linear representation of \(\mathfrak{g}_{\mathbf{C}}\) on a finite dimensional complex vector space V. Restricting to \(\mathfrak{g}\) , we obtain a representation of \(\mathfrak{g}\) on the real vector space \(V_{(\mathbf{R})}\) ; since G is simply-connected, there exists an analytic linear representation \(\pi\) of G on \(V_{(\mathbf{R})}\) such that \(\rho = L(\pi)\) . It follows from Prop. 7 of no. 3 that the image \(\delta(X(T))\) of X(T) in \(\mathfrak{t}_{\mathbf{C}}^{*}\) contains all the weights of \(\rho\) on V. This being true for every representation \(\rho\) of \(\mathfrak{g}_{\mathbf{C}}\) , it follows from Chap. VIII, §7, no. 2, Th. 1 that \(\delta(X(T))\) contains the group of weights of \(\delta(R)\) , which is by definition the set of \(\lambda \in \mathfrak{t}_{\mathbf{C}}^{*}\) such that \(\lambda(H_{\delta(\alpha)}) \in \mathbf{Z}\) for all \(\alpha \in R\) , in other words, \(\lambda(K_{\alpha}) \in 2\pi i\mathbf{Z}\) for all \(\alpha \in R\) . Thus, the group X(T) contains all the elements \(\lambda\) of X(T) \(\otimes\) Q such that \(\langle\lambda, K_{\alpha} \rangle \in \mathbf{Z}\) for all \(\alpha \in R\) , which implies by duality that \(\Gamma(T)\) is generated by the \(K_{\alpha}\) , hence the assertion \(A(G,T)\) .
\item
  G is the direct product of a simply-connected group \(G'\) and a torus S. Then T is the direct product of a maximal torus \(T'\) of \(G'\) with S, \(\varGamma(\mathrm{T})\) can be identified with \(\varGamma(\mathrm{T}')\times\varGamma(\mathrm{S})\) , \(\pi_{1}(\mathrm{G})\) with \(\pi_{1}(\mathrm{G}')\times\pi_{1}(\mathrm{S})\) , and \(f(\mathrm{G},\mathrm{T})\) with the homomorphisms with components \(f(\mathrm{G}^{\prime},\mathrm{T}^{\prime})\) and \(f(\mathrm{S},\mathrm{S})\) . Since \(f(\mathrm{S},\mathrm{S})\) is bijective, the canonical map \(\varGamma(\mathrm{T}^{\prime})\to\varGamma(\mathrm{T})\) maps \(\operatorname{Ker}f(\mathrm{G}^{\prime},\mathrm{T}^{\prime})\) bijectively onto \(\operatorname{Ker}f(\mathrm{G},\mathrm{T})\) . Moreover, the \(K_{\alpha}\) belong to the Lie algebra of the derived group \(G^{\prime}\) of G, hence to the image of \(\varGamma(\mathrm{T}^{\prime})\) , so it is immediate that \(A(\mathrm{G}^{\prime},\mathrm{T}^{\prime})\) implies \(A(\mathrm{G},\mathrm{T})\) , hence assertion \(A(\mathrm{G},\mathrm{T})\) , in view of a).
\item
  General case. There exists a surjective morphism \(p : G' \to G\) with finite kernel, where \(G'\) is the direct product of a simply-connected group by a torus ( \(\S1\) , no. 4, Prop. 4). If \(T'\) is the inverse image of T in \(G'\) (this is a maximal torus of \(G'\) by \(\S2\) , no. 3, Prop. 1), and N the kernel of p, we have exact sequences \(0 \to N \to G' \to G \to 0\) and \(0 \to N \to T' \to T \to 0\) , hence a commutative diagram with exact rows (no. 2, Remark 1 and General Topology, Chap. XI, in preparation)
\end{enumerate}

\[
\begin{array}{c c c c c c c c c} 0 & \longrightarrow & \Gamma (\mathrm{T} ^ {\prime}) & \longrightarrow & \Gamma (\mathrm{T}) & \longrightarrow & \mathrm{N} & \longrightarrow & 0 \\ & & \Big \downarrow f (\mathrm{G} ^ {\prime}, \mathrm{T} ^ {\prime}) & & \Big \downarrow f (\mathrm{G}, \mathrm{T}) & & \Big \downarrow \mathrm{Id} _ {\mathrm{N}} \\ 0 & \longrightarrow & \pi_ {1} (\mathrm{G} ^ {\prime}) & \longrightarrow & \pi_ {1} (\mathrm{G}) & \longrightarrow & \mathrm{N} & \longrightarrow & 0. \end{array}
\]

It follows immediately from the snake diagram (Algebra, Chap. X, p.~4, Prop. 2) that \(A(\mathrm{G}', \mathrm{T}')\) implies \(A(\mathrm{G}, \mathrm{T})\) , hence the proposition, in view of \(b\) ).

COROLLARY 1. G is simply-connected if and only if the family \((K_{\alpha})_{\alpha \in \mathrm{R}(\mathrm{G},\mathrm{T})}\) generates \(\Gamma (\mathrm{T})\) .

COROLLARY 2. Let H be a connected closed subgroup of G containing T; there is an exact sequence

\[
0 \longrightarrow \mathrm{N} (\mathrm{H}, \mathrm{T}) \longrightarrow \mathrm{N} (\mathrm{G}, \mathrm{T}) \longrightarrow \pi_ {1} (\mathrm{H}) \longrightarrow \pi_ {1} (\mathrm{G}) \longrightarrow 0.
\]

This follows from Algebra, Chap. X, p.~4, Prop. 2 (snake diagram), applied to the commutative diagram

\[
\begin{array}{c c c c c c c c c} 0 & \longrightarrow & \mathrm{N(H,T)} & \longrightarrow & \varGamma (\mathrm{T}) & \longrightarrow & \pi_ {1} (\mathrm{H}) & \longrightarrow & 0 \\ & & \Big \downarrow & & \Big \downarrow & & \Big \downarrow & & \\ 0 & \longrightarrow & \mathrm{N(G,T)} & \longrightarrow & \varGamma (\mathrm{T}) & \longrightarrow & \pi_ {1} (\mathrm{G}) & \longrightarrow & 0. \end{array}
\]

Remark. It can be shown (cf.~Exercise 2 of §5) that \(\pi_2(\mathrm{G / H})\) is zero. The exactness of the preceding sequence then gives an isomorphism from \(\pi_2(\mathrm{G / H})\) to \(\mathrm{N(G,T) / N(H,T)}\) .

COROLLARY 3. The homomorphism \(\pi_1(\mathrm{D}(\mathrm{G})) \to \pi_1(\mathrm{G})\) corresponding to the inclusion of \(\mathrm{D}(\mathrm{G})\) into \(\mathrm{G}\) induces an isomorphism from \(\pi_1(\mathrm{D}(\mathrm{G}))\) to the torsion subgroup of \(\pi_1(\mathrm{G})\) .

Indeed, \(\mathrm{T} \cap \mathrm{D}(\mathrm{G})\) is a maximal torus of \(\mathrm{D}(\mathrm{G})\) ( \(\S 2\) , no. 3, Prop. 1 \(c\) ); from the exact sequence

\[
0 \longrightarrow \Gamma (\mathrm{T} \cap \mathrm{D} (\mathrm{G})) \longrightarrow \Gamma (\mathrm{T}) \longrightarrow \Gamma (\mathrm{T} / (\mathrm{T} \cap \mathrm{D} (\mathrm{G}))) \longrightarrow 0
\]

and Proposition 11, we obtain an exact sequence

\[
0 \longrightarrow \pi_ {1} (\mathrm{D} (\mathrm{G})) \longrightarrow \pi_ {1} (\mathrm{G}) \longrightarrow \Gamma (\mathrm{T} / (\mathrm{T} \cap \mathrm{D} (\mathrm{G}))) \longrightarrow 0,
\]

hence the corollary, since \(\pi_{1}(\mathrm{D}(\mathrm{G}))\) is finite and \(\Gamma(\mathrm{T}/(\mathrm{T}\cap\mathrm{D}(\mathrm{G})))\) is free.

